\textbf{(CAP-COLUNI 2021 -- Q13)} Em fevereiro de 2020, um supermercado vendia uma mercadoria A por R\$13,00. No mês de março, em virtude da Pandemia da Covid-19, o Brasil entrou em quarentena e a economia do país foi afetada. Assim, em maio de 2020, a mercadoria A passou a custar R\$16,51, no mesmo supermercado.

Considerando o sistema de juros simples, podemos afirmar que a taxa de juros ao mês é dada por:

\begin{enumerate}
	\item 27\% a.m.
	\item 1,17\% a.m.
	\item 11,7\% a.m.
	\item 9\% a.m.
\end{enumerate}
